\documentclass[10pt,a4paper]{amsart}
\usepackage{amsmath,amssymb,amsthm,graphicx}
% Retained preamble for the Stacks Project chapter {CH} (shared project preamble).
% The chapter file pinned-original-source.tex only carries \input{{preamble}} on its first line,
% so no preamble text lives inside the chapter itself.  The Stacks Project uses one single shared
% preamble for every chapter: the snapshot stages it three times (sources/stacks-chow/preamble.tex,
% sources/stacks-etale-cohomology/preamble.tex and focused-seeds/AGP/sources/stacks-snapshot-20261011/
% preamble.tex) and those copies are identical in size.  The block below is copied mechanically,
% line for line, from sources/stacks-chow/preamble.tex lines 14-18 and 154-217: the Xy-pic setup and
% the native theorem declarations together with the mathematical macros.  Nothing was added, removed
% or redefined; the wrapper class and the amsmath/amssymb/amsthm/graphicx base packages are supplied
% by the assembler.
% For commutative diagrams we use Xy-pic
% Chapter: stacks-spaces-more-morphisms.tex -- More on Morphisms of Spaces.
\usepackage[all]{xy}

% We use 2cell for 2-commutative diagrams.
\xyoption{2cell}

% Theorem environments.
%
\theoremstyle{plain}
\newtheorem{theorem}[subsection]{Theorem}
\newtheorem{proposition}[subsection]{Proposition}
\newtheorem{lemma}[subsection]{Lemma}

\theoremstyle{definition}
\newtheorem{definition}[subsection]{Definition}
\newtheorem{example}[subsection]{Example}
\newtheorem{exercise}[subsection]{Exercise}
\newtheorem{situation}[subsection]{Situation}

\theoremstyle{remark}
\newtheorem{remark}[subsection]{Remark}
\newtheorem{remarks}[subsection]{Remarks}

\numberwithin{equation}{subsection}

% Macros
%
\def\lim{\mathop{\mathrm{lim}}\nolimits}
\def\colim{\mathop{\mathrm{colim}}\nolimits}
\def\Spec{\mathop{\mathrm{Spec}}}
\def\Hom{\mathop{\mathrm{Hom}}\nolimits}
\def\Ext{\mathop{\mathrm{Ext}}\nolimits}
\def\SheafHom{\mathop{\mathcal{H}\!\mathit{om}}\nolimits}
\def\SheafExt{\mathop{\mathcal{E}\!\mathit{xt}}\nolimits}
\def\Sch{\mathit{Sch}}
\def\Mor{\mathop{\mathrm{Mor}}\nolimits}
\def\Ob{\mathop{\mathrm{Ob}}\nolimits}
\def\Sh{\mathop{\mathit{Sh}}\nolimits}
\def\NL{\mathop{N\!L}\nolimits}
\def\CH{\mathop{\mathrm{CH}}\nolimits}
\def\proetale{{pro\text{-}\acute{e}tale}}
\def\etale{{\acute{e}tale}}
\def\QCoh{\mathit{QCoh}}
\def\Ker{\mathop{\mathrm{Ker}}}
\def\Im{\mathop{\mathrm{Im}}}
\def\Coker{\mathop{\mathrm{Coker}}}
\def\Coim{\mathop{\mathrm{Coim}}}

% Boxtimes
%
\DeclareMathSymbol{\boxtimes}{\mathbin}{AMSa}{"02}

%
% Macros for moduli stacks/spaces
%
\def\QCohstack{\mathcal{QC}\!\mathit{oh}}
\def\Cohstack{\mathcal{C}\!\mathit{oh}}
\def\Spacesstack{\mathcal{S}\!\mathit{paces}}
\def\Quotfunctor{\mathrm{Quot}}
\def\Hilbfunctor{\mathrm{Hilb}}
\def\Curvesstack{\mathcal{C}\!\mathit{urves}}
\def\Polarizedstack{\mathcal{P}\!\mathit{olarized}}
\def\Complexesstack{\mathcal{C}\!\mathit{omplexes}}
% \Pic is the operator that assigns to X its picard group, usage \Pic(X)
% \Picardstack_{X/B} denotes the Picard stack of X over B
% \Picardfunctor_{X/B} denotes the Picard functor of X over B
\def\Pic{\mathop{\mathrm{Pic}}\nolimits}
\def\Picardstack{\mathcal{P}\!\mathit{ic}}
\def\Picardfunctor{\mathrm{Pic}}
\def\Deformationcategory{\mathcal{D}\!\mathit{ef}}
\newenvironment{reference}{}{}
\newenvironment{slogan}{}{}
\newenvironment{history}{}{}
\usepackage{cleveref}
\crefname{originallemma}{original lemma}{original lemmas}
\Crefname{originallemma}{Original lemma}{Original lemmas}
\makeatletter
\newlabel{etale-theorem-remarkable-equivalence}{{etale-theorem-remarkable-equivalence}{source}{}{sourcepointer-1}{}}
\newlabel{etale-theorem-remarkable-equivalence@cref}{{[originallemma][0][]etale-theorem-remarkable-equivalence}{[1][0][]source}{}{}{}}
\makeatother
\crefname{originallemma}{original lemma}{original lemmas}
\Crefname{originallemma}{Original lemma}{Original lemmas}
\makeatletter
\newlabel{lemma-composition-formally-smooth-etale-unramified}{{lemma-composition-formally-smooth-etale-unramified}{source}{}{sourcepointer-2}{}}
\newlabel{lemma-composition-formally-smooth-etale-unramified@cref}{{[originallemma][0][]lemma-composition-formally-smooth-etale-unramified}{[1][0][]source}{}{}{}}
\makeatother
\crefname{originallemma}{original lemma}{original lemmas}
\Crefname{originallemma}{Original lemma}{Original lemmas}
\makeatletter
\newlabel{lemma-representable-property-formally-property}{{lemma-representable-property-formally-property}{source}{}{sourcepointer-3}{}}
\newlabel{lemma-representable-property-formally-property@cref}{{[originallemma][0][]lemma-representable-property-formally-property}{[1][0][]source}{}{}{}}
\makeatother
\title{An etale morphism on top of a formally unramified morphism stays formally unramified}
\author{The Stacks Project Authors}
\date{stacks-project:spaces-more-morphisms.tex:local-snapshot-20261011; GNU FDL 1.2 or later}
\begin{document}
\maketitle
\noindent\textit{Editorial note.} Statement span is lines 2566-2602 of sources/stacks-spaces-more-morphisms/pinned-original-source.tex; proof span is lines 2604-2654, the complete author proof (49 body lines).\par\medskip
\noindent\textit{Editorial note.} Navigation stubs: 3 cited-but-not-included labels are supplied as non-load-bearing pointers with their pinned tags where the snapshot tag index has them.\par\medskip
\noindent\textit{Editorial note.} extra\_packages is empty on purpose: the retained preamble already loads xy with options and a second declaration clashes in the assembler wrapper.\par\medskip
\noindent\textit{Editorial note.} Text is retained verbatim, including the author's own typos and placeholders.\par\medskip
\noindent\textit{Allowed standard prerequisites (editorial).}\begin{itemize}
\item Explicit prerequisite boundary. Formally unramified morphisms and their lifting criterion, and the lifting property of etale morphisms; the argument is retained.
\end{itemize}\medskip
\ifdefined\hypertarget\hypertarget{sourcepointer-1}{}\fi
\noindent\textit{Source reference.} Original lemma etale-theorem-remarkable-equivalence: sources/stacks-spaces-more-morphisms/pinned-original-source.tex, source lines 2654--2654. Cited by the retained statement or proof but not part of the collected spans; supplied only as a navigation pointer. The argument does not depend on its content.\par\medskip
\ifdefined\hypertarget\hypertarget{sourcepointer-2}{}\fi
\noindent\textit{Source reference.} Original lemma lemma-composition-formally-smooth-etale-unramified: sources/stacks-spaces-more-morphisms/pinned-original-source.tex, source lines 2486--2495. Cited by the retained statement or proof but not part of the collected spans; supplied only as a navigation pointer. The argument does not depend on its content.\par\medskip
\ifdefined\hypertarget\hypertarget{sourcepointer-3}{}\fi
\noindent\textit{Source reference.} Original lemma lemma-representable-property-formally-property: sources/stacks-spaces-more-morphisms/pinned-original-source.tex, source lines 2529--2537. Cited by the retained statement or proof but not part of the collected spans; supplied only as a navigation pointer. The argument does not depend on its content.\par\medskip
\begin{lemma}
\label{lemma-etale-on-top}
Let $S$ be a scheme contained in $\Sch_{fppf}$.
Let $F, G, H : (\Sch/S)_{fppf}^{opp} \to \textit{Sets}$.
Let $a : F \to G$, $b : G \to H$ be transformations of functors.
Assume that $a$ is representable, surjective, and \'etale.
\begin{enumerate}
\item If $b$ is formally smooth, then $b \circ a$ is formally smooth.
\item If $b$ is formally \'etale, then $b \circ a$ is formally \'etale.
\item If $b$ is formally unramified, then $b \circ a$ is formally unramified.
\end{enumerate}
Conversely, consider a solid commutative diagram
$$
\xymatrix{
G \ar[d]_b & T \ar[d]^i \ar[l] \\
H & T' \ar[l] \ar@{-->}[lu]
}
$$
with $T'$ an affine scheme over $S$
and $i : T \to T'$ a closed immersion defined by an ideal of square zero.
\begin{enumerate}
\item[(4)] If $b \circ a$ is formally smooth, then for every $t \in T$
there exists an \'etale morphism of affines $U' \to T'$ and a morphism
$U' \to G$ such that
$$
\xymatrix{
G \ar[d]_b & T \ar[l] & T \times_{T'} U' \ar[d] \ar[l]\\
H & T' \ar[l] & U' \ar[llu] \ar[l]
}
$$
commutes and $t$ is in the image of $U' \to T'$.
\item[(5)] If $b \circ a$ is formally unramified, then there exists at most
one dotted arrow in the diagram above, i.e., $b$ is formally unramified.
\item[(6)] If $b \circ a$ is formally \'etale, then there exists exactly one
dotted arrow in the diagram above, i.e., $b$ is formally \'etale.
\end{enumerate}
\end{lemma}

\begin{proof}
Assume $b$ is formally smooth (resp.\ formally \'etale,
resp.\ formally unramified). Since an \'etale morphism is both
smooth and unramified we see that $a$ is representable and smooth
(resp.\ \'etale, resp. unramified). Hence parts (1), (2) and (3)
follow from a combination of
Lemma \ref{lemma-representable-property-formally-property}
and
Lemma \ref{lemma-composition-formally-smooth-etale-unramified}.

\medskip\noindent
Assume that $b \circ a$ is formally smooth. Consider a diagram
as in the statement of the lemma. Let $W = F \times_G T$.
By assumption $W$ is a scheme surjective \'etale over $T$. By
\'Etale Morphisms, Theorem \ref{etale-theorem-remarkable-equivalence}
there exists a scheme $W'$ \'etale over $T'$ such that $W = T \times_{T'} W'$.
Choose an affine open subscheme $U' \subset W'$ such that $t$ is in
the image of $U' \to T'$. Because $b \circ a$ is formally
smooth we see that the exist morphisms $U' \to F$ such that
$$
\xymatrix{
F \ar[d]_{b \circ a} & W \ar[l] & T \times_{T'} U' \ar[d] \ar[l]\\
H & T' \ar[l] & U' \ar[llu] \ar[l]
}
$$
commutes. Taking the composition $U' \to F \to G$ gives a
map as in part (5) of the lemma.

\medskip\noindent
Assume that $f, g : T' \to G$ are two dotted arrows fitting into the
diagram of the lemma. Let $W = F \times_G T$.
By assumption $W$ is a scheme surjective \'etale over $T$. By
\'Etale Morphisms, Theorem \ref{etale-theorem-remarkable-equivalence}
there exists a scheme $W'$ \'etale over $T'$ such that $W = T \times_{T'} W'$.
Since $a$ is formally \'etale the compositions
$$
W' \to T' \xrightarrow{f} G
\quad\text{and}\quad
W' \to T' \xrightarrow{g} G
$$
lift to morphisms $f', g' : W' \to F$ (lift on affine opens and glue by
uniqueness). Now if $b \circ a : F \to H$ is formally unramified, then
$f' = g'$ and hence $f = g$ as $W' \to T'$ is an \'etale covering. This proves
part (6) of the lemma.

\medskip\noindent
Assume that $b \circ a$ is formally \'etale. Then by part (4) we
can \'etale locally on $T'$ find a dotted arrow fitting into the diagram
and by part (5) this dotted arrow is unique. Hence we may glue the
local solutions to get assertion (6). Some details omitted.
\end{proof}

\end{document}
